\section{Làm quen giao diện}
\label{sec:giao-dien}

Dưới đây là màn hình chính bạn sẽ thấy khi mở một dự án, với các con số đánh
dấu từng khu vực chức năng quan trọng.

\anhcochuthich{\linewidth}{giao-dien}{
  \chuthich{0.10}{0.89}{1}
  \chuthich{0.013}{0.70}{2}
  \chuthich{0.36}{0.62}{3}
  \chuthich{0.77}{0.55}{4}
  \chuthich{0.88}{0.90}{5}
  \chuthich{0.90}{0.985}{6}
  \chuthich{0.48}{0.985}{7}
  \chuthich{0.53}{0.90}{8}
  \chuthich{0.86}{0.06}{9}
  \chuthich{0.987}{0.86}{10}
}{Màn hình soạn thảo của SumiTeX}

\begin{enumerate}
  \item \textbf{Tệp} (Files): Hiển thị toàn bộ tệp trong dự án; bạn chỉ cần
    nhấp vào tên tệp để mở.
  \item \textbf{Thanh bên trái}: Chứa các công cụ quản lý như Tệp, Mục lục, Tìm
    kiếm, Lịch sử và Tích hợp. Để thu gọn bảng đang mở, bạn hãy nhấp lại vào
    chính biểu tượng đó.
  \item \textbf{Trình soạn thảo}: Khu vực chính để bạn nhập mã và viết nội
    dung.
  \item \textbf{Bản PDF}: Hiển thị kết quả sau khi biên dịch. Khi nhấp đúp vào
    một chữ bất kỳ trên PDF, con trỏ sẽ tự động nhảy tới dòng mã tương ứng ở
    bên trái.
  \item \nut{Code} / \nut{Visual}: Chuyển đổi giữa chế độ xem mã nguồn
    \LaTeX{} và chế độ hiển thị trực quan gần giống bản in.
  \item \textbf{Nút chạy} (▶): Biên dịch tài liệu sang PDF (hoặc dùng phím tắt
    \phim{Ctrl}+\phim{S}).
  \item \textbf{Ô tìm kiếm}: Mở nhanh \emph{bảng lệnh} (xem chi tiết tại
    chương~\ref{sec:bang-lenh}).
  \item \textbf{Đèn báo lỗi}: Cho biết số lượng lỗi và cảnh báo hiện có trong
    tệp đang mở.
  \item \nut{Download}: Tải bản PDF hoàn chỉnh về máy tính.
  \item \textbf{Thanh bên phải}: Nơi xem nhận xét, thông báo và trao đổi trực
    tiếp với những người cùng soạn thảo.
\end{enumerate}

\begin{meo}
Bạn có thể thay đổi nhanh bố cục làm việc bằng các phím tắt
\phim{Ctrl}+\phim{1} (chỉ hiện mã), \phim{Ctrl}+\phim{2} (hiện song song mã và
PDF) hoặc \phim{Ctrl}+\phim{3} (chỉ hiện PDF). Ngoài ra, khi muốn tập trung tối
đa vào việc viết, hãy nhấn \phim{Ctrl}+\phim{K} rồi \phim{Z} để vào chế độ
\emph{Zen}, và nhấn \phim{Esc} để thoát.
\end{meo}

\begin{thuchanh}
Trong sân tập, bạn hãy mở tệp \texttt{bai-01.tex}, sau đó thử chuyển qua lại
giữa ba bố cục bằng \phim{Ctrl}+\phim{1}, \phim{Ctrl}+\phim{2} và
\phim{Ctrl}+\phim{3}.
\end{thuchanh}